% 2010 年浙江卷文科第 14 题：各行、各列都成等差数列的数表。
% 来源 PDF：普通高考/2010/2010浙江文.pdf，SHA-256 a73e22ea405697775b2d5cd82532fcffad0f506b849d9f199265b4297b8a754b。
% 原页：第 1 页（A3 三栏）第 2 栏下部，书页码 641。
% 重绘依据：照原页全框格排版，表头一行为“第 1 列、第 2 列、第 3 列、⋯”，
%   行标为“第 1 行、第 2 行、第 3 行、⋯”，格内数字 1 2 3 / 2 4 6 / 3 6 9，末行末列为省略号。
\begin{tikzpicture}[scale=0.9, semithick, every node/.style={inner sep=1pt, font=\small}]

  % 竖线
  \foreach \x in {0, 1.95, 3.5, 5.05, 6.6, 7.7} {
    \draw (\x,0.66) -- (\x,-2.64);
  }
  % 横线
  \foreach \y in {0.66, 0, -0.66, -1.32, -1.98, -2.64} {
    \draw (0,\y) -- (7.7,\y);
  }

  % 表头
  \node at (0.975,0.33) {第 $1$ 列};
  \node at (2.725,0.33) {第 $2$ 列};
  \node at (4.275,0.33) {第 $3$ 列};
  \node at (6.05,0.33) {$\cdots$};

  % 第 1 行
  \node at (0.975,-0.33) {第 $1$ 行};
  \node at (2.725,-0.33) {1};
  \node at (4.275,-0.33) {2};
  \node at (5.825,-0.33) {3};
  \node at (7.15,-0.33) {$\cdots$};

  % 第 2 行
  \node at (0.975,-0.99) {第 $2$ 行};
  \node at (2.725,-0.99) {2};
  \node at (4.275,-0.99) {4};
  \node at (5.825,-0.99) {6};
  \node at (7.15,-0.99) {$\cdots$};

  % 第 3 行
  \node at (0.975,-1.65) {第 $3$ 行};
  \node at (2.725,-1.65) {3};
  \node at (4.275,-1.65) {6};
  \node at (5.825,-1.65) {9};
  \node at (7.15,-1.65) {$\cdots$};

  % 末行
  \node at (0.975,-2.31) {$\cdots$};
  \node at (2.725,-2.31) {$\cdots$};
  \node at (4.275,-2.31) {$\cdots$};
  \node at (5.825,-2.31) {$\cdots$};
  \node at (7.15,-2.31) {$\cdots$};
\end{tikzpicture}
